% !TEX root = ../../ctfp-print.tex

\lettrine[lhang=0.17]{我}{们见过幺半群的好几种表述}: 作为一个集合、作为单物件范畴、作为幺半范畴里
的一个物件. 这个简单的概念里还能再榨出多少油水?

来试一试. 把幺半群定义成一个集合 $m$ 加上一对函数:
\begin{align*}
  \mu  & \Colon m\times{}m \to m \\
  \eta & \Colon 1 \to m
\end{align*}
这里 $1$ 是 $\Set$ 里的终止物件 —— 单元素集合. 第一个函数定义乘法 (它取一对元素, 返回它们
的积), 第二个从 $m$ 里挑出单位元. 并非随意为这两个函数挑了这种签名就得到一个幺半群, 为此还
得施加额外的条件: 结合律和单位律. 不过暂且把这些放到一边, 只考虑 ``潜在幺半群''. 一对函数是
两个函数集合的笛卡尔积里的一个元素. 我们知道, 这两个集合可以表示成指数物件:
\begin{align*}
  \mu  & \in m^{m\times{}m} \\
  \eta & \in m^1
\end{align*}
这两个集合的笛卡尔积是:
\[m^{m\times{}m}\times{}m^1\]
动用一点中学代数 (它在每个笛卡尔闭范畴里都行得通), 可以把它改写成:
\[m^{m\times{}m + 1}\]
其中的 $+$ 代表 $\Set$ 里的余积. 我们刚刚用单个函数替换掉了一对函数 —— 它是这个函数集合
的一个元素:
\[m\times{}m + 1 \to m\]
这个函数集合的任意一个元素都是一个潜在幺半群.

这种表述的美妙之处在于, 它能引出一些有趣的一般化. 比如说, 用这套语言我们怎么描述群? 群就是
幺半群再加一个函数, 它给每个元素指派自己的逆元, 类型是 $m \to m$. 举个例子, 整数就构成一个
群, 二元运算是加法, 单位元是零, 逆元是取负. 要定义一个群, 我们得从一组三个函数出发:
\begin{align*}
  m\times{}m \to m \\
  m \to m          \\
  1 \to m
\end{align*}
和先前一样, 我们可以把这一组三个函数合进一个函数集合:
\[m\times{}m + m + 1 \to m\]
我们开头有一个二元运算符 (加法)、一个一元运算符 (取负)、一个零元运算符 (单位元 —— 这里是
零), 把它们合成了一个函数. 所有这种签名的函数都定义了潜在的群.

我们还可以这样一路走下去. 比如要定义环, 就再加一个二元运算符和一个零元运算符, 如此等等. 每
一次我们最后得到的都是一个函数类型, 它的左边是一串幂的和 (可能包含零次幂 —— 也就是终止物
件), 右边是集合本身.

现在我们可以痛痛快快地推广了. 头一桩, 把集合换成物件、把函数换成态射. 我们把 n 元运算符定义
成从 n 元积出发的态射, 这意味着需要一个支持有限积的范畴. 对零元运算符, 我们要求终止物件存在,
所以要一个笛卡尔范畴. 为了把这些运算符合起来, 又需要指数, 于是这得是个笛卡尔闭范畴. 最后,
我们还需要余积来给这通代数戏法收尾.

或者, 我们干脆忘掉这些公式是怎么推出来的, 只盯着最后那个积. 我们那个态射左边的积之和定义了
一个自函子. 要是我们索性挑一个任意的自函子 $F$ 呢? 那样一来, 我们对范畴就不必施加任何约束
了. 我们得到的东西叫作 F-代数.

F-代数是一个三元组, 包含一个自函子 $F$、一个物件 $a$, 以及一个态射
\[F\ a \to a\]
这个物件通常叫作载体、一个底层的物件, 或者在编程的语境里叫作载体 \emph{类型}. 这个态射通
常叫作求值函数或者结构映射. 不妨把函子 $F$ 想成是在生成表达式, 把那个态射想成是在求值.

下面是 F-代数的 Haskell 定义:

\src{snippet01}
它把一个代数与它的求值函数等同了起来.

在幺半群那个例子里, 我们说的函子是:

\src{snippet02}
这就是 $1 + a\times{}a$ 的 Haskell 写法 (还记得
\hyperref[simple-algebraic-data-types]{代数数据类型}吗).

环可以用下面这个函子来定义:

\src{snippet03}
它是 $1 + 1 + a\times{}a + a\times{}a + a$ 的 Haskell 写法.

整数组就是环的一个例子. 我们可以选 \code{Integer} 当载体类型, 把求值函数定义成:

\src{snippet04}
基于同一个函子 \code{RingF} 还有更多 F-代数. 比如多项式构成一个环, 方阵也构成一个环.

如你所见, 函子的角色是生成表达式, 再由代数的求值器把它们算出来. 到目前为止我们见过的表达式
都很简单. 我们感兴趣的往往是更精细复杂的表达式, 它们可以用递归定义出来.

\section{递归}

生成任意表达式树的一个办法, 是把函子定义里的变量 \code{a} 换成递归. 举例来说, 环里的任意
表达式都由这个树状数据结构生成:

\src{snippet05}
我们可以把最初的环求值器换成它的递归版本:

\src{snippet06}
这仍然不太实用, 因为我们被迫把所有整数都表示成一堆 1 的加和, 不过急起来也够用了.

可我们该怎么用 F-代数的语言描述表达式树? 我们得想办法把这样一件事形式化: 在我们函子的定义里,
把那个自由的类型变量递归地替换成替换的结果. 设想一步一步来做. 首先把深度为 1 的树定义成:

\src{snippet07}
我们是在用 \code{RingF a} 生成的深度为 0 的树, 去填 \code{RingF} 定义里的空洞. 深度为 2
的树照同样的办法得到:

\src{snippet08}
它也可以写成:

\src{snippet09}
继续这个过程, 我们可以写下一个符号方程:

\begin{snipv}
type RingF\textsubscript{n+1} a = RingF (RingF\textsubscript{n} a)
\end{snipv}
从概念上讲, 把这个过程重复无穷多次, 我们最后就得到了自己的 \code{Expr}. 注意 \code{Expr}
并不依赖 \code{a}. 我们从哪儿起步无所谓, 总会落在同一个地方. 对任意范畴里的任意自函子来说
这并不总是成立, 但在范畴 $\Set$ 里事情很美好.

当然了, 这是个手挥式的论证, 我稍后会把它弄严格.

无穷多次施加一个自函子会产生一个 \newterm{fixed point 不动点}, 它是一个这样定义的物件:
\[Fix\ f = f\ (Fix\ f)\]
这个定义背后的直觉是, 既然我们已经施加了 $f$ 无穷多次才得到 $Fix\ f$, 再多施加一次也
不会有任何改变. 在 Haskell 里, 不动点的定义是:

\src{snippet10}
可以说, 要是构造器的名字跟它所定义的类型的名字不同, 这段会更好读, 就像这样:

\src{snippet11}
不过我还是沿用通行的记法. 构造器 \code{Fix} (你要是乐意, 也可以叫它 \code{In}) 可以被看成
一个函数:

\src{snippet12}
还有另一个函数, 它剥掉一层函子应用:

\src{snippet13}
这两个函数互为逆. 我们稍后会用到它们.

\section{F-代数构成的范畴}

这是书里最老的一招: 每当你想出一个构造新物件的办法, 就看看它们能不能构成一个范畴. 毫不奇怪,
给定自函子 $F$ 之上的代数确实构成一个范畴. 那个范畴里的物件是代数 —— 由一个载体物件 $a$
和一个态射 $F\ a \to a$ 组成的对, 两者都来自原来的范畴 $\cat{C}$.

要把图景补齐, 我们还得定义 F-代数范畴里的态射. 一个态射必须把一个代数 $(a, f)$ 映到另一个
代数 $(b, g)$. 我们把它定义成一个在载体之间做映射的态射 $m$ —— 在原来的范畴里它从 $a$
到 $b$. 随便一个态射可不行: 我们要求它同两个求值器相容. (我们把这样保持结构的态射叫作
\newterm{homomorphism 同态}.) F-代数的同态是这么定义的. 首先注意, 我们可以把 $m$ 提升到
这个映射:
\[F\ m \Colon F\ a \to F\ b\]
接着在它后面接上 $g$, 就能到达 $b$. 等价地, 我们可以先用 $f$ 从 $F\ a$ 走到 $a$,
再接上 $m$. 我们要求这两条路径相等:
\[g \circ F\ m = m \circ f\]

\begin{figure}[H]
  \centering
  \includegraphics[width=0.3\textwidth]{images/alg.png}
\end{figure}

\noindent
很容易就能说服自己这确实是一个范畴 (提示: $\cat{C}$ 的恒等态射就挺好, 而同态的复合还是
一个同态).

F-代数范畴里的初始物件 (如果它存在的话) 叫作 \newterm{initial algebra 初始代数}. 我们把这个
初始代数的载体叫作 $i$, 它的求值器叫作 $j \Colon F\ i \to i$. 结果发现, 初始代数的求值器
$j$ 是一个同构. 这个结论叫作 Lambek 定理. 证明依赖初始物件的定义, 它要求从初始代数出发,
到任何别的 F-代数都存在唯一的同态 $m$. 既然 $m$ 是一个同态, 下图就必须交换:

\begin{figure}[H]
  \centering
  \includegraphics[width=0.3\textwidth]{images/alg2.png}
\end{figure}

\noindent
现在来构造一个以 $F\ i$ 为载体的代数. 这样一个代数的求值器必须是从 $F\ (F\ i)$ 到
$F\ i$ 的态射. 我们只要把 $j$ 提升一下, 就能轻易造出这样的求值器:
\[F\ j \Colon F\ (F\ i) \to F\ i\]
因为 $(i, j)$ 是初始代数, 所以从它到 $(F\ i, F\ j)$ 必定存在唯一的同态 $m$. 下图必须
交换:

\begin{figure}[H]
  \centering
  \includegraphics[width=0.3\textwidth]{images/alg3a.png}
\end{figure}

\noindent
但我们还有这个平凡交换的图 (两条路径压根是同一个!):

\begin{figure}[H]
  \centering
  \includegraphics[width=0.3\textwidth]{images/alg3.png}
\end{figure}

\noindent
它可以解读成: $j$ 是一个代数之间的同态, 把 $(F\ i, F\ j)$ 映到 $(i, j)$. 我们可以把
这两张图粘在一起, 得到:

\begin{figure}[H]
  \centering
  \includegraphics[width=0.6\textwidth]{images/alg4.png}
\end{figure}

\noindent
这张图反过来又可以解读成: $j \circ m$ 是一个代数同态, 只不过这回两个代数是同一个. 而且,
因为 $(i, j)$ 是初始的, 从它到自己只能有一个同态, 那就是恒等态射 $\id_i$ —— 我们知道
它就是一个代数同态. 于是 $j \circ m = \id_i$. 借助这个事实和左边那张图的交换性质, 我们
可以证明 $m \circ j = \id_{Fi}$. 这说明 $m$ 是 $j$ 的逆, 从而 $j$ 是 $F\ i$ 与 $i$
之间的同构:
\[F\ i \cong i\]
而这正是在说 $i$ 是 $F$ 的一个不动点. 这就是当初那个手挥式论证背后的形式证明.

回到 Haskell: 我们认出 $i$ 就是我们的 \code{Fix f}, $j$ 就是我们的构造器 \code{Fix},
它的逆就是 \code{unFix}. Lambek 定理里的那个同构告诉我们, 要得到初始代数, 我们只要拿函子
$f$, 把它的参数 $a$ 换成 \code{Fix f}. 我们也由此明白了不动点为什么不依赖 $a$.

\section{自然数}

自然数也可以定义成一个 F-代数. 出发点是一对态射:
\begin{align*}
  zero & \Colon 1 \to N \\
  succ & \Colon N \to N
\end{align*}
第一个挑出零, 第二个把所有数映到它们的后继. 和先前一样, 我们可以把这两个合成一个:
\[1 + N \to N\]
左边定义了一个函子, 用 Haskell 可以写成这样:

\src{snippet14}
这个函子的不动点 (它生成的初始代数) 在 Haskell 里可以编码成:

\src{snippet15}
一个自然数要么是零, 要么是另一个数的后继. 这就是所谓自然数的皮亚诺表示.

\section{折叠态射}

我们用 Haskell 记法把初始性条件重写一遍. 初始代数我们叫它 \code{Fix f}, 它的求值器就是构造器
\code{Fix}. 从初始代数出发, 到同一个函子之上的任何别的代数, 都存在唯一的态射 \code{m}.
我们挑一个以 \code{a} 为载体、以求值器为 \code{alg} 的代数.

\begin{figure}[H]
  \centering
  \includegraphics[width=0.4\textwidth]{images/alg5.png}
\end{figure}

\noindent
顺带请注意 \code{m} 是什么: 它是不动点的求值器, 是整棵递归表达式树的求值器. 我们来找一个
实现它的一般办法.

Lambek 定理告诉我们, 构造器 \code{Fix} 是一个同构, 我们把它的逆叫作 \code{unFix}. 因此
我们可以把这张图里的一个箭头翻过来, 得到:

\begin{figure}[H]
  \centering
  \includegraphics[width=0.4\textwidth]{images/alg6.png}
\end{figure}

\noindent
把这张图的交换条件写下来:

\begin{snip}{haskell}
m = alg . fmap m . unFix
\end{snip}
我们可以把这个方程解读成 \code{m} 的一个递归定义. 对任何用函子 \code{f} 造出的有限树,
这个递归都注定会终止. 这一点只要看到 \code{fmap m} 是在函子 \code{f} 的顶层之下起作用的
就明白了. 换句话说, 它作用在原树的子节点上, 而子节点永远比原树浅一层.

我们把 \code{m} 施加到一棵用 \code{Fix\ f} 构造出来的树上时, 会发生这样的事:
\code{unFix} 的动作剥掉构造器, 露出树的顶层; 然后我们把 \code{m} 施加到顶层节点的每一个
孩子上, 产出 \code{a} 类型的结果; 最后我们施加那个非递归的求值器 \code{alg}, 把这些结果
合起来. 关键在于, 我们的求值器 \code{alg} 是一个简单的、不带递归的函数.

既然对任何代数 \code{alg} 我们都能这么干, 那么定义一个高阶函数就很合理: 它把代数当作参数,
交出我们称为 \code{m} 的那个函数. 这个高阶函数叫作折叠态射 (catamorphism):

\src{snippet16}
看一个例子. 取定义自然数的那个函子:

\src{snippet17}
我们选 \code{(Int, Int)} 当载体类型, 并把代数定义成:

\src{snippet18}
你很容易就能说服自己: 这个代数的折叠态射 \code{cata fib} 算的是斐波那契数.

一般说来, \code{NatF} 的一个代数定义了一个递推关系: 当前元素的值由前一个元素决定. 折叠态射
则把那个序列的第 n 个元素算出来.

\section{折叠}

\code{e} 的列表是下面这个函子的初始代数:

\src{snippet19}
事实上, 把变量 \code{a} 换成递归的结果 (我们叫它 \code{List e}), 我们得到:

\src{snippet20}
列表函子的一个代数挑定一个具体的载体类型, 并定义一个对两个构造器做模式匹配的函数. 它在
\code{NilF} 上的值告诉我们要如何求值一个空列表, 它在 \code{ConsF} 上的值告诉我们如何把
当前元素与先前累积的值合起来.

举例来说, 下面这个代数可以用来算列表的长度 (载体类型是 \code{Int}):

\src{snippet21}
的确, 由此得到的折叠态射 \code{cata lenAlg} 算的就是列表的长度. 注意这个求值器是两样东西
的组合: (1) 一个函数, 它取一个列表元素和一个累加值, 返回新的累加值; (2) 一个起始值, 这里
是零. 值的类型和累加器的类型都由载体类型给出.

把它同 Haskell 的传统定义比较一下:

\src{snippet22}
\code{foldr} 的两个参数恰好就是这个代数的两个分量.

再试一个例子:

\src{snippet23}
同样, 把它与这个比较一下:

\src{snippet24}
如你所见, \code{foldr} 不过是折叠态射在列表上的一种方便的专门化.

\section{余代数}

照例, 我们有 F-余代数的对偶构造, 其中态射的方向反了过来:
\[a \to F\ a\]
给定函子的余代数同样构成一个范畴, 其中的同态保持余代数结构. 那个范畴里的终止物件
$(t, u)$ 叫作终止 (或最终) 余代数. 对每一个别的代数 $(a, f)$, 都存在唯一的同态 $m$,
让下图交换:

\begin{figure}[H]
  \centering
  \includegraphics[width=0.4\textwidth]{images/alg7.png}
\end{figure}

\noindent
终止余代数是函子的一个不动点, 意思是态射 $u \Colon t \to F\ t$ 是一个同构 (余代数版的
Lambek 定理):
\[F\ t \cong t\]
在编程里, 终止余代数通常被解读成一份生成 (可能无限的) 数据结构或者迁移系统的配方.

正如折叠态射可以用来对一个初始代数求值, 展开态射 (anamorphism) 可以用来对一个终止余代数
做余求值:

\src{snippet25}
余代数的一个典型例子基于这样一个函子, 它的不动点是 \code{e} 类型的元素构成的无穷流. 这个
函子是:

\src{snippet26}
它的不动点是:

\src{snippet27}
\code{StreamF e} 的一个余代数是一个函数, 它取一个 \code{a} 类型的种子, 产出一个对
(\code{StreamF} 是个花哨的名字, 说的就是对), 由一个元素和下一个种子组成.

你很容易就能造出产生无穷序列的简单余代数, 比如平方数列表, 或者倒数列表.

更有意思一点的例子是产生素数列表的余代数. 诀窍是用一个无穷列表当载体. 我们的起始种子是列表
\code{{[}2..{]}}. 下一个种子是这个列表删掉所有 2 的倍数之后剩下的尾部, 也就是从 3 开始的
奇数列表. 下一步我们取这个列表的尾部, 删掉所有 3 的倍数, 如此下去. 你大概已经认出这是埃拉托
斯特尼筛法的架势. 这个余代数由下面这个函数实现:

\src{snippet28}
这个余代数的展开态射生成了素数列表:

\src{snippet29}
流是一个无穷列表, 所以把它转成 Haskell 列表应该是可行的. 为此我们可以用同一个函子
\code{StreamF} 构造一个代数, 然后在它上面跑一个折叠态射. 比如, 下面就是一个把流转成列表的
折叠态射:

\src{snippet30}
这里, 同一个不动点同时是同一个自函子的初始代数和终止余代数. 在任意范畴里事情并不总是这样.
一般说来, 一个自函子可以有很多 (或者一个也没有) 不动点. 初始代数就是所谓的最小不动点,
终止余代数则是最大不动点. 不过在 Haskell 里, 两者由同一个公式定义, 它们重合了.

列表的展开态射叫作 unfold. 要造出有限的列表, 人们把函子改成产出一个 \code{Maybe} 对:

\src{snippet31}
\code{Nothing} 的值会终止列表的生成.

余代数的一个有意思的例子同透镜有关. 一个透镜可以表示成一对取值器与置值器:

\src{snippet32}
这里 \code{a} 通常是某个带 \code{s} 类型字段的积数据类型. 取值器取出那个字段的值, 置值器
则把这个字段换成新值. 这两个函数可以合成一个:

\src{snippet33}
我们还可以把这个函数进一步改写成:

\src{snippet34}
其中我们定义了一个函子:

\src{snippet35}
注意这可不是一个由积之和搭起来的简单代数函子, 它牵扯到一个指数 $a^s$.

透镜就是以 \code{a} 为载体类型的、这个函子的一个余代数. 我们先前见过, \code{Store s}
同时也是一个余单子. 结果发现, 一个表现良好的透镜对应于一个与余单子结构相容的余代数. 下一节
我们再来谈这件事.

\section{挑战}

\begin{enumerate}
  \tightlist
  \item
        实现一元多项式环的求值函数. 你可以把一个多项式表示成 $x$ 的各次幂前面的系数列表.
        举例来说, $4x^2-1$ 可以表示成 (从零次幂写起) \code{{[}-1, 0, 4{]}}.
  \item
        把上一个构造推广到多个独立变量的多项式, 比如 $x^2y-3y^3z$.
  \item
        实现 $2\times{}2$ 矩阵环的那个代数.
  \item
        定义一个余代数, 让它的展开态射产生自然数的平方列表.
  \item
        用 \code{unfoldr} 生成前 $n$ 个素数的列表.
\end{enumerate}
